\apendice{Documentación de usuario}

\section{Introducción}

El presente apéndice constituye la documentación orientada al usuario final del
sistema HIVES (\textit{Honey Identification and Verification by Spectroscopy}).
Su objetivo es proporcionar las instrucciones necesarias para la correcta
instalación, configuración y operación de la aplicación, sin requerir
conocimientos de programación por parte del usuario.

El sistema ha sido diseñado para ser operado por dos perfiles de usuario
diferenciados:

\begin{itemize}
    \item \textbf{Usuario técnico o semitécnico:} personal de laboratorio,
    técnicos de calidad apícola o investigadores con conocimientos básicos
    en el uso de un ordenador. Este perfil es capaz de gestionar la
    configuración del puerto serie y la calibración del sensor de forma
    autónoma.
    \item \textbf{Usuario no técnico:} apicultores u operarios sin
    conocimientos informáticos avanzados, cuya interacción con el sistema
    se limita al flujo principal de análisis. Para este perfil se recomienda
    que la configuración inicial sea realizada por personal técnico.
\end{itemize}

La interfaz gráfica sigue principios de diseño minimalista, minimizando la
curva de aprendizaje y permitiendo obtener resultados de clasificación de
forma intuitiva y eficiente.
\imagentamano{img/Interfaz/MainView_hives.png}{Vista general de la ventana principal de la aplicación HIVES. Fuente: Elaboración Propia.}{1}

\section{Requisitos de usuario}

Para garantizar el correcto funcionamiento del sistema HIVES, el equipo sobre
el que se ejecute la aplicación deberá cumplir los siguientes requisitos mínimos
de hardware y software según el sistema operativo empleado:
\newpage
\tablaSmallSinColores
  {Requisitos mínimos del sistema para la ejecución de HIVES.}
  {p{3.5cm} p{8cm}}
  {requisitos_sistema}
  {
    \textbf{Componente} & \textbf{Especificación mínima} \\
  }
  {
    Sistema operativo  & Windows 10 (64\,bits) o superior; Ubuntu 20.04 LTS
                         o superior; macOS 11 (\textit{Big Sur}) o superior \\
    Procesador         & Intel Core i3 (o equivalente AMD), 2{,}0\,GHz,
                         o ARM (M1, M2,...) en el caso de macOS \\
    Memoria RAM        & 4\,GB \\
    Almacenamiento     & 2\,GB de espacio libre en disco \\
    Puerto de conexión & Puerto USB libre (para adaptador USB--Serie) \\
    Conectividad       & No requerida (operación completamente local) \\
  }

Adicionalmente, es necesario que el espectrómetro se encuentre conectado al
equipo mediante el adaptador USB--Serie correspondiente y que el sistema
operativo haya reconocido correctamente el dispositivo antes de iniciar
la aplicación.

\section{Instalación}

La distribución de HIVES se realiza en forma de paquete autocontenido generado
mediante la herramienta PyInstaller en formato \emph{one-dir}. Dicho formato
encapsula el intérprete de Python, todas las dependencias y el modelo de
Inteligencia Artificial, eximiendo al usuario de cualquier instalación técnica
previa. El resultado es una carpeta de aplicación (\texttt{HIVES/}) que contiene
el ejecutable principal junto a un subdirectorio \texttt{\_internal/} con las
bibliotecas, los binarios de soporte y el modelo entrenado. Se proporciona un
paquete independiente para cada sistema operativo compatible.

\begin{quote}
\textbf{Importante:} la carpeta de aplicación debe copiarse y conservarse
\emph{íntegra}. El ejecutable principal depende del subdirectorio
\texttt{\_internal/} que lo acompaña; si se extrae el ejecutable de forma aislada,
la aplicación no podrá iniciarse. En macOS esta restricción no aplica, ya que la
distribución adopta la forma de un \emph{bundle} \texttt{.app} autocontenido.
\end{quote}

\subsection*{Windows}
\begin{enumerate}
  \item Copiar la carpeta completa \texttt{HIVES\textbackslash} en el directorio de
        trabajo deseado (por ejemplo,
        \texttt{C:\textbackslash Usuarios\textbackslash <usuario>\textbackslash HIVES\textbackslash}),
        conservando en su interior el ejecutable \texttt{HIVES.exe} y el
        subdirectorio \texttt{\_internal\textbackslash}.
  \item Conectar el espectrómetro al ordenador mediante el cable USB antes de
        iniciar la aplicación. El sensor se comunica con el equipo a través de la
        placa Arduino Uno que lo gobierna. Verificar en el \emph{Administrador de
        dispositivos} que el dispositivo aparece reconocido y que se le ha asignado
        un puerto COM activo (por ejemplo, \texttt{COM3}).
  \item Hacer doble clic sobre \texttt{HIVES.exe} para iniciar la aplicación. No se
        requiere instalación adicional ni permisos de administrador.
  \item Asegurarse de que la carpeta donde reside la aplicación cuenta con permisos
        de escritura, ya que la base de datos local y los reportes exportados se
        almacenan en esa misma ubicación.
\end{enumerate}

\subsection*{Linux}
\begin{enumerate}
  \item Copiar la carpeta completa \texttt{HIVES/} en el directorio de trabajo
        deseado, conservando en su interior el binario \texttt{HIVES} y el
        subdirectorio \texttt{\_internal/}.
  \item Conectar el espectrómetro al ordenador mediante el cable USB. La placa
        Arduino que gobierna el sensor se enumera en Linux con un nombre de
        dispositivo que depende de su circuito de comunicación: las placas con
        conversor USB--serie (CH340, FTDI o CP210x) aparecen como
        \texttt{/dev/ttyUSB0}, mientras que las placas con USB nativo lo hacen como
        \texttt{/dev/ttyACM0}. Para conocer el puerto asignado, listar ambos
        conjuntos de dispositivos:
\begin{verbatim}
ls /dev/ttyUSB* /dev/ttyACM* 2>/dev/null
\end{verbatim}
        Si hubiera varios dispositivos conectados y existiera duda sobre cuál
        corresponde al espectrómetro, puede desconectarse y volverse a conectar la
        placa y consultar los últimos mensajes del núcleo, que indican el nombre
        exacto recién asignado:
\begin{verbatim}
dmesg | grep -i tty
\end{verbatim}
  \item Conceder al usuario permiso de acceso al puerto serie. Los dispositivos
        \texttt{ttyUSB} y \texttt{ttyACM} pertenecen al grupo \texttt{dialout}, por
        lo que basta con añadir el usuario a dicho grupo (solución persistente y
        recomendada) y conceder permiso de ejecución al binario:
\begin{verbatim}
chmod +x HIVES/HIVES
sudo usermod -aG dialout $USER
\end{verbatim}
        Es necesario cerrar la sesión y volver a iniciarla para que el cambio de
        grupo surta efecto. Como alternativa puntual para una única sesión, puede
        otorgarse acceso directo al dispositivo concreto (sustituyendo el nombre por
        el identificado en el paso anterior):
\begin{verbatim}
sudo chmod a+rw /dev/ttyUSB0
\end{verbatim}
  \item Ejecutar la aplicación desde el terminal, situándose en la carpeta de la
        aplicación, o haciendo doble clic sobre el binario si el entorno de
        escritorio lo permite:
\begin{verbatim}
cd HIVES && ./HIVES
\end{verbatim}
\end{enumerate}

\subsection*{macOS (Apple Silicon, M1 o superior)}
\begin{enumerate}
  \item Copiar el paquete \texttt{HIVES.app} en la carpeta \emph{Aplicaciones} o en
        el directorio de trabajo deseado.

        \textbf{Nota:} el ejecutable distribuido ha sido compilado para la
        arquitectura ARM64 (Apple Silicon). No ha sido probado en otras arquitecturas. Para 
        equipos Mac con procesador Intel, se requerirá el uso de Rosetta~2.
  \item Conectar el espectrómetro al ordenador mediante el cable USB. El sistema
        operativo reconocerá automáticamente la placa Arduino Uno y le asignará un
        puerto con un nombre similar a \texttt{/dev/cu.usbmodem1401}, lo que puede
        verificarse desde la terminal con:
\begin{verbatim}
ls /dev/cu.*
\end{verbatim}
  \item En caso de que macOS bloquee la ejecución por provenir de un desarrollador
        no identificado, acceder a \emph{Ajustes del Sistema} $\rightarrow$
        \emph{Privacidad y seguridad} y autorizar manualmente la ejecución de la
        aplicación.
  \item Abrir la aplicación haciendo doble clic sobre \texttt{HIVES.app}.
\end{enumerate}

\section{Manual del usuario}

Esta sección describe el flujo de trabajo completo de la aplicación HIVES,
detallando las pantallas disponibles y las acciones que el usuario puede
realizar en cada una de ellas.

\subsection*{Estructura de la interfaz}

La aplicación está organizada en tres elementos principales:

\begin{itemize}
    \item \textbf{Menú lateral izquierdo:} panel de control desde el que se
    gestiona la conexión con el Arduino Uno, se selecciona el puerto serie
    activo y se accede a los parámetros de configuración del análisis.
    \item \textbf{Ventana principal:} área central de operación donde se
    inicia el análisis, se visualizan los resultados y se exportan los reportes.
    \item \textbf{Ventana de históricos:} repositorio de consulta de todos los
    análisis previos registrados, con capacidad de filtrado por fecha,
    categoría y nivel de confianza. Además se pueden extraer los datos a un fichero \textbf{CSV}.
    \item \textbf{Ventana de muestras:} repositorio de consulta de todas las muestras tomadas hasta la fecha, con capacidad de filtrado por fecha, y de extracción de las mismas en \textbf{CSV}.
\end{itemize}

\subsection*{Paso 1 — Configuración inicial (perfil técnico)}

\textit{Este paso está destinado al usuario técnico. El usuario no técnico
puede omitirlo si la configuración ya ha sido realizada previamente.}

\begin{enumerate}
    \item En el \textbf{menú lateral izquierdo} de la ventana principal,
    localizar el selector de puerto serie.
    \item Seleccionar el puerto al que está conectado el Arduino Uno. La lista
    se actualiza automáticamente con los puertos disponibles detectados por
    el sistema operativo.
    \item Ajustar los parámetros de análisis adicionales si fuera necesario.
    La aplicación conservará esta configuración en sesiones sucesivas.
\end{enumerate}

\imagentamano{img/Interfaz/BerMenu_hives.png}{Menú lateral izquierdo de la ventana principal: selector de puerto serie y parámetros de análisis. Fuente: Elaboración Propia.}{0.2}

\subsection*{Paso 2 — Calibración del sensor (perfil técnico)}

Se recomienda realizar una calibración del sensor al inicio de cada sesión
de trabajo, especialmente si el dispositivo ha sido desconectado o si las
condiciones ambientales han variado respecto a la sesión anterior.

\begin{enumerate}
    \item Desde la ventana principal, accionar el control de
    \textbf{Calibración}.
    \item Seguir las indicaciones que aparezcan en pantalla durante el proceso. Habrá que hacer 2 calibraciones: una de blancos y otra de negros. Para ello se usará un papel de reflectancia en el primer caso, y una cubeta negra en el segundo caso.
    \item Una vez completada la calibración, el sistema estará listo para
    realizar análisis con garantías de precisión.
\end{enumerate}

\subsection*{Paso 3 — Análisis de una muestra}

Este es el flujo principal de uso de la aplicación, accesible para cualquier
perfil de usuario.

\begin{enumerate}
    \item Colocar la muestra de miel en el soporte del espectrómetro conforme
    al procedimiento físico establecido.
    \item En la \textbf{ventana principal}, accionar el control
    \textbf{«Iniciar Análisis»}.
    \item El sistema establecerá la conexión con el espectrómetro, capturará
    los datos espectrales y ejecutará el modelo de clasificación de forma
    automática. Durante este proceso se mostrará un indicador de progreso en
    pantalla.
    \item Transcurrido el tiempo de análisis (máximo 2~minutos), los resultados
    se visualizarán en la interfaz: \textbf{categoría de miel clasificada},
    \textbf{índice de confianza del modelo} y \textbf{métricas espectrales}
    de la muestra.
    \item Los datos quedan registrados automáticamente en la base de datos
    local sin intervención del usuario.
\end{enumerate}
\imagentamano{img/Interfaz/DataCapture.png}{Ventana principal mostrando los resultados tras completar un análisis de muestra. Fuente: Elaboración Propia.}{}


\subsection*{Paso 4 — Exportación del reporte PDF}

\begin{enumerate}
    \item Accionar el control \textbf{«Generar Reporte PDF»} en la ventana
    principal.
    \item El sistema abrirá un cuadro de diálogo del sistema operativo para
    seleccionar la ruta y el nombre del archivo de destino.
    \item Confirmar la ruta. El documento PDF se generará y guardará en el
    directorio indicado, quedando disponible para su consulta, impresión o
    distribución externa.
\end{enumerate}

\subsection*{Paso 5 — Consulta del histórico (perfil técnico)}

\begin{enumerate}
    \item Accionar el control \textbf{«Ver Histórico»} en la ventana principal.
    \item La \textbf{ventana de históricos} mostrará una tabla con todos los
    análisis registrados: fecha, tipo de miel clasificado y nivel de confianza.
    \item Utilizar los controles de filtrado disponibles (rango de fechas,
    categoría de miel, umbral de confianza) para acotar los resultados.
    \item Seleccionar un registro de la tabla para visualizar su detalle
    completo o iniciar la exportación PDF descrita en el paso anterior.
\end{enumerate}

\imagentamano{img/Interfaz/HistoricView.png}{Ventana de históricos con los filtros de búsqueda disponibles. Fuente: Elaboración Propia.}{1}
\newpage
\subsection*{Mensajes de error frecuentes}
\tablaSmallSinColores
  {Mensajes de error de la aplicación y acciones correctoras.}
  {p{4.5cm} p{8cm}}
  {mensajes_error}
  {
    \textbf{Mensaje} & \textbf{Acción recomendada} \\
  }
  {
    «No se detecta el espectrómetro» &
      Verificar que el cable USB del Arduino Uno está correctamente
      conectado. Comprobar que el puerto serie seleccionado en la
      ventana de configuración coincide con el asignado por el sistema
      operativo (\texttt{COMx} en Windows, \texttt{/dev/ttyACM0} en
      Linux, \texttt{/dev/cu.usbmodem*} en macOS). En Linux, verificar
      además que el usuario pertenece al grupo \texttt{dialout}. \\
    «Error al cargar el modelo de IA» &
     Comprobar que el archivo del modelo (\texttt{mejor\_modelo.keras}) y el fichero
    de clases (\texttt{clases.json}) se encuentran en la carpeta \texttt{models/}
    junto al ejecutable. \\
    «Sin permisos de escritura» &
      Mover el ejecutable a un directorio con permisos de escritura
      para el usuario actual. En macOS y Linux, verificar los permisos
      del directorio con \texttt{ls -la}. \\
    «No se han encontrado coincidencias» &
      Ampliar los criterios de filtrado en la ventana de históricos. \\
  }